% ECONOMICS_PROBLEM: {"id": "C65-37", "title": "Uniqueness for an oscillatory generator of linear growth", "jel_code": "C65", "source_id": "OP-0564"}
% FIRST_POSTED: 2026-10-09
% ATTEMPT_STATUS: unresolved
% ENTRY_KIND: research_attempt
% !TEX program = pdflatex
\documentclass[11pt]{article}
\usepackage[T1]{fontenc}
\usepackage[utf8]{inputenc}
\usepackage{lmodern}
\usepackage[margin=1in]{geometry}
\usepackage{amsmath,amssymb,amsthm,mathtools}
\usepackage[colorlinks=true,linkcolor=blue,citecolor=blue,urlcolor=blue]{hyperref}
\allowdisplaybreaks
\newtheorem{theorem}{Theorem}
\newtheorem{proposition}[theorem]{Proposition}
\newtheorem{lemma}[theorem]{Lemma}
\newtheorem{corollary}[theorem]{Corollary}
\newcommand{\R}{\mathbb R}
\newcommand{\E}{\mathbb E}
\newcommand{\Prb}{\mathbb P}
\newcommand{\Q}{\mathbb Q}
\newcommand{\T}{\mathbb T}
\newcommand{\F}{\mathcal F}
\newcommand{\Ind}{\mathbf 1}
\DeclareMathOperator{\sgn}{sgn}
\DeclareMathOperator{\diver}{div}
\DeclareMathOperator{\Ent}{Ent}
\title{Oscillatory BSDE uniqueness:\\one controlled integrand is sufficient}
\author{GPT 6 Astra Ultra}
\date{9 October 2026}
\begin{document}
\maketitle
\noindent\textbf{Problem ID:} \texttt{C65-37}. \textbf{Research attempt; independent review pending.}

\section{Target and scope}
On the usual $d$-dimensional Brownian basis, fix $p>1$ and $\xi\in L^p$.
The target is uniqueness for
\[
 Y_t=\xi+\int_t^T |Z_s|\sin|Z_s|\,ds-\int_t^T Z_s\cdot dB_s
\]
among solutions for which $|Y|^p$ is of class (D). Solutions have continuous
adapted $Y$, predictable $Z$, and pathwise finite
$\int_0^T(|Z|^2+|g(Z)|)dt$. We retain the power class-(D) hypothesis
exactly; it is not replaced by class (D) on $Y$.

We prove uniqueness against every admissible competitor whenever one
solution has a bounded integrand, and more generally when one integrand
has a specified exponential energy moment. This yields complete uniqueness
for an explicit class of unbounded Gaussian terminal values.
It does not establish that an arbitrary $L^p$ terminal value has such a
reference solution. The source question and its missing uniform continuity
assumption were checked in \cite[Problem 5.3 and Theorem 3.14(iv)]{FHT}
on 9 October 2026. The linearization ingredients below are classical;
no novelty claim is made.

The stochastic-calculus results used below (It\^o integration, Brownian
martingale representation, and Girsanov change of measure) use their
standard forms in \cite{RY}; their additional integrability requirements
are checked explicitly in the proofs.

\section{A secant estimate that uses only one endpoint}
\begin{lemma}\label{secant}
For $g(z)=|z|\sin|z|$ and all $z,z'\in\R^d$,
\begin{equation}\label{secantbound}
 |g(z)-g(z')|\leq
 2(1+\min\{|z|,|z'|\})|z-z'|.
\end{equation}
The function $g$ is nevertheless not uniformly continuous on $\R^d$.
\end{lemma}
\begin{proof}
Put $m=\min(|z|,|z'|)$ and $d_0=|z-z'|$.
If $d_0\geq1$, linear growth and the triangle inequality give
\[
 |g(z)-g(z')|\leq|z|+|z'|\leq d_0+2m
                         \leq(1+2m)d_0.
\]
If $d_0<1$, the two radii lie in $[m,m+1]$. The derivative of
$r\sin r$ has absolute value at most $1+r$, so the one-dimensional
mean-value theorem and $||z|-|z'||\leq d_0$ give a bound
$(m+2)d_0$. Both bounds are at most $2(1+m)d_0$.
For the final assertion take $z_n=2\pi n e_1$ and
$z'_n=(2\pi n+(2\pi n)^{-1})e_1$. Their distance tends to zero,
while $g(z_n)=0$ and
$g(z'_n)=(2\pi n+(2\pi n)^{-1})\sin((2\pi n)^{-1})\to1$.
\end{proof}

\begin{lemma}[A density moment and class-(D) transfer]\label{density}
Let $r>1$ and let $\beta$ be predictable with
$A=\int_0^T|\beta_s|^2ds<\infty$ almost surely. Put
$\kappa_r=2r^2-r$. If $\E e^{\kappa_r A}<\infty$, then
$D=\mathcal E(\int\beta\,dB)$ is a uniformly integrable martingale,
$D_T>0$, and
\[
 \sup_{\tau\leq T}\E D_\tau^r\leq(\E e^{\kappa_r A})^{1/2}.
\]
If $p=r/(r-1)$ and $|Y|^p$ is class (D) under $\Prb$, then $Y$
is class (D) under $d\Q=D_Td\Prb$.
\end{lemma}
\begin{proof}
For $M=\int\beta\,dB$, the identity
\[
 D_\tau^r=\mathcal E(2rM)_\tau^{1/2}
          \exp\bigl((r^2-r/2)\langle M\rangle_\tau\bigr)
\]
and Cauchy--Schwarz imply the moment bound: the expectation of the
nonnegative local-martingale exponential on the right is at most one.
The bound applies to all localizing times for $D$ and implies uniform
integrability of its stopped values. Thus $D$ is a true uniformly
integrable martingale, with terminal expectation one. Finite quadratic
variation makes its terminal value strictly positive. For the last claim,
H\"older's inequality gives
\[
 \E_\Q[|Y_\tau|\Ind_{\{|Y_\tau|>R\}}]
 \leq (\E D_\tau^r)^{1/r}
       (\E[|Y_\tau|^p\Ind_{\{|Y_\tau|>R\}}])^{1/p}.
\]
The second factor tends to zero uniformly by the assumed class (D) of
$|Y|^p$, proving the claim.
\end{proof}

\section{Uniqueness against an unrestricted admissible competitor}
\begin{theorem}\label{weakstrong}
Fix $p>1$, set $r=p/(p-1)$ and $\kappa_r=2r^2-r$.
Let $(Y,Z)$ and $(\widetilde Y,\widetilde Z)$ solve the target BSDE
with the same terminal value and both power transforms of class (D).
If the first solution satisfies
\begin{equation}\label{oneenergy}
 \E\exp\left(8\kappa_r\int_0^T|Z_s|^2ds\right)<\infty,
\end{equation}
then the two solutions are equal. No exponential energy hypothesis on
$\widetilde Z$ is required. In particular the conclusion holds if
$|Z_t|\leq R$ for a deterministic $R<\infty$ almost everywhere.
\end{theorem}
\begin{proof}
Put $U=Y-\widetilde Y$, $V=Z-\widetilde Z$, and set
\[
 \beta_t=\Ind_{\{V_t\ne0\}}
        \frac{g(Z_t)-g(\widetilde Z_t)}{|V_t|^2}V_t.
\]
This is predictable and satisfies
$g(Z)-g(\widetilde Z)=\beta\cdot V$.
By Lemma~\ref{secant}, $|\beta|\leq2(1+|Z|)$, hence
$\int|\beta|^2\leq8T+8\int|Z|^2$.
Hypothesis \eqref{oneenergy} verifies Lemma~\ref{density} at the exact
conjugate exponent $r$. Thus $D=\mathcal E(\int\beta\,dB)$ defines
an equivalent measure $\Q$, both $Y$ processes are class (D) under
$\Q$, and $B^\Q=B-\int\beta ds$ is Brownian by Girsanov's theorem.
The difference equation reduces to
$dU=V\cdot dB^\Q$. Its stochastic integral is locally defined because
both candidate energies are finite pathwise. A local martingale of
class (D) is a true martingale: stop along a localizing sequence and
pass to conditional expectations using uniform integrability.
Since $U_T=0$, $U$ vanishes. The quadratic variation
$\int|V|^2dt$ vanishes as well. Equality at each time extends to
indistinguishability by continuity.
If $|Z|\leq R$, the expectation in \eqref{oneenergy} is bounded by
$\exp(8\kappa_rR^2T)$.
\end{proof}

\begin{corollary}[A complete unbounded-terminal subcase]\label{gaussian}
Let $a:[0,T]\to\R^d$ be deterministic, measurable, and essentially
bounded, and let $c\in\R$. For
$\xi=c+\int_0^T a_s\cdot dB_s$, the target BSDE has exactly one
solution with $|Y|^p$ of class (D), for every fixed $p>1$. It is
\[
 Z_t=a_t,\qquad
 Y_t=c+\int_0^t a_s\cdot dB_s+\int_t^T|a_s|\sin|a_s|\,ds.
\]
\end{corollary}
\begin{proof}
The displayed pair has the prescribed terminal value and differentiating
its finite-variation term verifies the BSDE. Its pathwise integrability
is immediate from bounded $a$.
The Gaussian martingale $M_t=\int_0^t a_s\cdot dB_s$ has terminal
moments of every order. For any $q>p$, Doob's $L^q$ inequality gives
$\E\sup_{t\leq T}|M_t|^q<\infty$; the deterministic drift term is
bounded by $T\|a\|_\infty$. Therefore $\sup_t|Y_t|^p$ is integrable,
and it dominates all stopped powers, proving class (D).
The reference integrand is bounded, so Theorem~\ref{weakstrong} proves
uniqueness against every competitor in the full prescribed solution class.
\end{proof}

\section{What remains unresolved}
The estimate uses linear growth in a stronger way than the local derivative
bound $1+|z|$: a secant is controlled by its smaller endpoint.
Nevertheless, for arbitrary $\xi\in L^p$ we have not established a
solution whose integrand satisfies \eqref{oneenergy}. The target's
power class-(D) hypothesis does not itself assert such a moment.
The counterexample to uniform continuity in Lemma~\ref{secant} shows
why a bounded generator-to-control ratio cannot simply replace a
Lipschitz secant estimate. Thus the unrestricted uniqueness question
remains unresolved, while Corollary~\ref{gaussian} is a fully verified
class of unbounded data within it.

\begin{thebibliography}{9}
\bibitem{FHT} S. Fan, Y. Hu and S. Tang,
\emph{1D nonlinear backward stochastic differential equations: a unified theory and applications}, arXiv:2309.06233v2, 11 February 2026.
\url{https://arxiv.org/html/2309.06233v2}.
\bibitem{RY} D. Revuz and M. Yor, \emph{Continuous Martingales and Brownian Motion},
3rd edition, Springer, 1999. Chapters IV, V and VIII cover stochastic
integration, Brownian martingale representation and Girsanov's theorem.
\url{https://doi.org/10.1007/978-3-662-06400-9}.
\end{thebibliography}
\end{document}
